\documentclass[class=article,border=10pt, margin=3mm]{standalone}
\usepackage{verbatim}
\usepackage{textcomp}
\usepackage{tikz}
\usetikzlibrary{shapes,arrows}
\usepackage[americanresistors,americaninductors,RPvoltages]{circuitikz}
\usepackage{amsmath}


\begin{document}

\begin{circuitikz}[american, scale=1.2]
    
    % 1. Definimos las coordenadas para controlar la altura (Rhombus height)
    % Si quieres el rombo más alto, aumenta el 2.0 y disminuye el -2.0
    \coordinate (T) at (3, 1.5);  % Top
    \coordinate (B) at (3, -1.5); % Bottom
    \coordinate (L) at (1.5, 0);  % Left
    \coordinate (R) at (4.5, 0);  % Right

    % --- Fuente de CA ---
    \draw (0, 1.5) to[vsourcesin, l=$V_{gen}$] (0, -1.5);
    \draw (0, 1.5) -- (T);   
    \draw (0, -1.5) -- (B); 

    % --- Diodos ---
    \draw (L) to[D*] (T); 
    \draw (L) to[D*] (B); 
    \draw (T) to[D*] (R); 
    \draw (B) to[D*] (R); 

    % 2. Nodo de tierra separado a la izquierda
    % Dibujamos una línea pequeña a la izquierda (-- ++(-0.5,0)) y luego la tierra
    \draw (L) -- ++(-0.2, 0) node[ground] {};

    % --- Resistencia de Carga RL ---
    \draw (R) -- (6, 0)
        to[R, name=res] (6, -2.5) 
        node[ground] {};

    % 3. Posicionamiento manual de etiquetas (para que no se choquen)
    % Usamos los "anchors" de la resistencia (res.west, res.east)
    
    % RL a la derecha de la resistencia
    \node[anchor=west] at (5.0, -0.8) {$R_L$};
    
    % Valor a la izquierda de la resistencia
    \node[anchor=east] at (5.8, -1.25) {$4.7\,\text{k}\Omega$};
   
    
    % Vout con signos + y - bien alineados
    \node[anchor=west] at (6.2, -0.5) {$+$};
    \node[anchor=west] at (6.2, -1.25) {$V_{out}$};
    \node[anchor=west] at (6.2, -2.0) {$-$};

\end{circuitikz}

\end{document}